\chapter{Sheaf and Stack Hyperstructures}
\label{ch:sheaf-hyperstructures}

\section{Overview}

Chapters~\ref{ch:semantic-manifold}--\ref{ch:typed-hypergraphs}
established the geometric, variational, and combinatorial underpinnings of the
RSVP--Polyxan universe. The present chapter introduces the sheaf- and
stack-theoretic hyperstructure that binds these layers into a coherent,
descent-compatible system.

Typed Polyxan hypergraphs (Chapter~\ref{ch:typed-hypergraphs}) form the
combinatorial substrate; RSVP fields provide geometric and variational
structure (Chapters~\ref{ch:rsvp-fields}--\ref{ch:cognitive-lagrangian});
cognitive rewrites induce local transformations
(Chapter~\ref{ch:lsystem}).  
But to maintain global consistency under curvature, stratification, semantic
flows, and higher-order rewrite coherences, we require a \emph{semantic stack
architecture}.

This chapter constructs:
\begin{itemize}
  \item the sheaf of semantic configurations,
  \item the stack of hypergraphs and rewrite objects,
  \item gluing and descent conditions,
  \item fibrational semantics over the RSVP manifold,
  \item and the homotopy-coherent structure needed for cognitive dynamics.
\end{itemize}

The resulting framework prepares the formal integration of RSVP fields and
Polyxan hypergraphs in Chapter~\ref{ch:coupling}.

\section{The Semantic Configuration Sheaf}

Let $\mathscr{P}$ be the Polyxan atom sheaf (Chapter~\ref{ch:semantic-manifold})
and $\mathscr{H}$ the hypergraph sheaf (Chapter~\ref{ch:typed-hypergraphs}).

\begin{definition}[Semantic Configuration Sheaf]
The sheaf of semantic configurations $\mathscr{C}$ assigns to each open set
$U\subseteq\mathcal{M}$ the category
\[
\mathscr{C}(U)
=
\mathrm{Conf}(U)
:=
\Gamma(U,\mathscr{P}^{\otimes \ast})
\;\times\;
\mathscr{H}(U),
\]
consisting of pairs of local atom configurations and local typed hypergraphs.
\end{definition}

\subsection{Restriction maps}

For inclusions $V\subseteq U$, the restriction
\[
\rho_{UV}:\mathscr{C}(U)\to\mathscr{C}(V)
\]
is defined componentwise:
\[
\rho_{UV}(X,H) = (X|_V,\; H|_V).
\]

These satisfy the usual sheaf axioms.

\section{Stacks of Semantic Structures}

Configurations alone cannot handle equivalences induced by
\emph{rewrite rules}, \emph{gauge symmetries}, \emph{stratified
automorphisms}, or \emph{hypergraph isomorphisms}.  
For this we require a stack.

\subsection{Stack of semantic objects}

\begin{definition}[Semantic Stack]
Define a prestack $\mathcal{S}$ by associating to each open $U$ the groupoid
\[
\mathcal{S}(U)
:= 
\mathrm{Iso}\bigl(\mathscr{C}(U)\bigr),
\]
whose objects are semantic configurations and whose morphisms are local
isomorphisms combining:
\begin{itemize}
  \item Polyxan atom equivalences,
  \item hypergraph isomorphisms,
  \item semantic coordinate transformations,
  \item and rewrite-generated equivalences.
\end{itemize}
\end{definition}

\begin{proposition}
$\mathcal{S}$ is a stack with respect to the Grothendieck topology on
$\mathcal{M}$.
\end{proposition}

\begin{proof}[Sketch]
Descent for objects follows from sheaf gluing for $\mathscr{P}$ and $\mathscr{H}$.
Descent for morphisms requires naturality under restrictions and follows from
the fact that semantic coordinate transformations are diffeomorphisms of
$\mathcal{M}$.
\end{proof}

\section{Fibration Structure and Semantic Transport}

\subsection{Grothendieck fibration}

Let $\pi:\mathcal{S}\to\mathcal{M}$ send each object $(X,H)$ to the region
$\mathrm{supp}(X,H)$.

\begin{definition}[Semantic Fibration]
The map $\pi:\mathcal{S}\to\mathcal{M}$ is a Grothendieck fibration, meaning:
for any arrow $f:x\to y$ in $\mathcal{M}$ and any $(X,H)$ over $y$, there is a
Cartesian lift over $f$.
\end{definition}

\begin{remark}
This ensures semantic structures transport coherently under RSVP flows and
rewrite-induced transformations.
\end{remark}

\subsection{Covariant vs.\ contravariant transport}

The fibration supports:
\[
f^\ast : \mathcal{S}(U) \to \mathcal{S}(V) \quad \text{(pullback)},
\]
\[
f_\ast : \mathcal{S}(V) \to \mathcal{S}(U) \quad \text{(pushforward)},
\]
each compatible with sheaf morphisms and rewrite rules.

RSVP vector fields $\mathbf{v}$ induce a flow $\varphi_t$ on $\mathcal{M}$ and
hence a transport functor:
\[
\varphi_t^\ast : \mathcal{S} \to \mathcal{S}.
\]

\section{Stratification and Higher Stacks}

\subsection{Stratified stack structure}

Strata $X_\alpha\subseteq\mathcal{M}$ (Chapter~\ref{ch:semantic-manifold})
produce a corresponding filtration of the semantic stack:
\[
\mathcal{S}_{\le k}
=
\left\{
(X,H)\in\mathcal{S}
:\;
\mathrm{supp}(X,H)\subseteq \bigcup_{\alpha\le k} X_\alpha
\right\}.
\]

\subsection{Higher morphisms}

Rewrite rules generate equivalences that must satisfy coherence relations.
This introduces 2-morphisms and higher.

\begin{definition}[Semantic 2-Groupoid]
The 2-groupoid $\mathcal{S}_{(2)}(U)$ has:
\begin{itemize}
  \item objects: semantic configurations on $U$,
  \item 1-morphisms: sheaf and hypergraph isomorphisms,
  \item 2-morphisms: rewrite coherence diagrams.
\end{itemize}
\end{definition}

In full generality, $\mathcal{S}$ is a \emph{weak $(\infty,1)$-stack}.

\section{Descent, Gluing, and Rewriting}

\subsection{Local rewrites and descent}

Rewrite rules $\rho$ apply locally:
\[
\rho: \mathscr{P}(U)\to\mathscr{P}(U)^{\otimes k}.
\]

By naturality (Chapter~\ref{ch:lsystem}), rewrites commute with restriction:
\[
\rho|_V = \rho\circ\rho_{UV}.
\]

Thus, rewrite steps satisfy descent.

\subsection{Gluing under rewrite coherence}

Consider a cover $\{U_i\}$ of $U$.  
Suppose $X_i\in\mathcal{S}(U_i)$ satisfy the cocycle conditions:
\[
X_i|_{U_{ij}} \cong X_j|_{U_{ij}}
\]
with coherence isomorphisms satisfying higher cocycle conditions.

\begin{proposition}[Rewrite-Compatible Descent]
If each $X_i$ is stable under a rewrite rule $\rho$, then the glued object
$X\in\mathcal{S}(U)$ is also stable under $\rho$.
\end{proposition}

\begin{proof}[Sketch]
Rewrite naturality and the stack condition ensure that local invariance implies
global invariance under descent.
\end{proof}

\section{Hypergraph Bundles and Associated Stacks}

A typed hypergraph $H$ can be organized as a fiber bundle over $\mathcal{M}$.

\subsection{Hypergraph bundle}

\begin{definition}[Hypergraph Bundle]
A hypergraph bundle is a diagram
\[
H \to \mathcal{M}
\]
such that for each $x\in\mathcal{M}$, the fiber $H_x$ is a small hypergraph
determined by the local type structure.
\end{definition}

\subsection{Associated stack}

Define the associated stack:
\[
\mathcal{H} := \left[\mathscr{H} / \mathrm{Aut}(\mathcal{M}) \right].
\]

Sections of $\mathcal{H}$ correspond to equivalence classes of embedded
hypergraphs.

\section{Semantic Boundary Conditions}

\subsection{Boundary sheaves}

For open embeddings $j:U\hookrightarrow \mathcal{M}$, define the boundary
sheaf:
\[
\partial\mathscr{C}(U)
=
\ker\left(\mathscr{C}(U)\to\mathscr{C}(\overline{U})\right).
\]

These capture content that “lives on the edges” of semantic regions.

\subsection{Boundary hypergraph conditions}

A hypergraph satisfies structural boundary conditions if:
\[
\forall e\in E,\quad \operatorname{supp}(e) \cap \partial U=\emptyset,
\]
or satisfies reflective/absorptive conditions depending on semantic context.

\section{Homotopy-Coherent Rewrite Semantics}

Rewrite operations are not strictly associative; coherence laws must be imposed.

\subsection{Homotopy coherence}

Define $n$-fold rewrite operations
\[
\rho^{(n)}: \mathscr{P}\to\mathscr{P}^{\otimes m_n}
\]
and coherence homotopies
\[
h_n : \rho^{(n)} \Rightarrow \rho^{(n+1)}
\]
satisfying Stasheff-like relations.

\begin{definition}[Cognitive $A_\infty$-structure]
Rewrite rules endow $\mathscr{P}$ with the structure of an
$A_\infty$-algebra object internal to the semantic stack.
\end{definition}

\subsection{BV compatibility}

The BV Laplacian $\Delta$ (Chapter~\ref{ch:cognitive-lagrangian}) must satisfy:
\[
\Delta^2 = 0,\qquad
\Delta(\rho) = 0,
\]
ensuring homological consistency.

\section{Summary}

This chapter introduced the sheaf and stack framework necessary for the
RSVP–Polyxan hyperstructure:
\begin{itemize}
  \item the sheaf $\mathscr{C}$ of semantic configurations,
  \item the semantic stack $\mathcal{S}$ capturing local-to-global structure,
  \item fibration and transport under RSVP flows,
  \item descent and gluing for rewrite rules,
  \item hypergraph bundles and associated stacks,
  \item boundary semantics and stratified structure,
  \item and higher homotopy-coherent rewrite semantics.
\end{itemize}

These constructs provide the formal machinery required to integrate RSVP fields
with typed Polyxan hypergraphs in the next chapter.
